
Reinforcement learning from execution feedback (RLEF) with binary rewards suffers from gradient starvation: at group size G=4, up to 69\% of GRPO training steps produce zero gradient because the model either passes or fails every completion in a group. This paper proposes offline variance-guided selection --- profiling a frozen model once to compute per-problem binary execution reward variance p\_i($1-p_{i})$, then selecting problems with the highest variance before training begins. On MBPP (374 training problems) with DeepSeek-Coder-7B-Instruct, two results are confirmed: (1) 7.8\% of problems exhibit nonzero variance under constrained generation (k=4, max\_new\_tokens=128), and (2) top-50 variance-guided selection achieves 4.$24\times$ higher mean variance than random-50 selection (Mann-Whitney U p=2.$22\times 10^{-6})$, stochastically dominating random selection across the full distribution, with profiling completing in approximately 32 seconds via vLLM. However, GRPO training on the variance-selected subset produces no gradient signal across more than 50,000 generation attempts (50--200 steps), because the model generates zero correct solutions under training conditions --- a cold-start regime that renders data selection method irrelevant regardless of analytical validity. Generation parameter mismatch between profiling (max\_new\_tokens=128, vLLM) and training (max\_completion\_length=512, HF GRPOTrainer) is identified as the primary candidate root cause among several alternatives. A prescriptive resolution protocol is proposed but not empirically validated. The primary contribution is a precisely characterized failure mode and its necessary precondition: cold-start verification is established as a prerequisite for any RLEF data selection method, and the first systematic per-problem empirical characterization of binary execution reward variance distribution on MBPP for a 7B code LLM is provided.

